% Options for packages loaded elsewhere
\PassOptionsToPackage{unicode}{hyperref}
\PassOptionsToPackage{hyphens}{url}
%
\documentclass[
]{book}
\usepackage{amsmath,amssymb}
\usepackage{iftex}
\ifPDFTeX
  \usepackage[T1]{fontenc}
  \usepackage[utf8]{inputenc}
  \usepackage{textcomp} % provide euro and other symbols
\else % if luatex or xetex
  \usepackage{unicode-math} % this also loads fontspec
  \defaultfontfeatures{Scale=MatchLowercase}
  \defaultfontfeatures[\rmfamily]{Ligatures=TeX,Scale=1}
\fi
\usepackage{lmodern}
\ifPDFTeX\else
  % xetex/luatex font selection
\fi
% Use upquote if available, for straight quotes in verbatim environments
\IfFileExists{upquote.sty}{\usepackage{upquote}}{}
\IfFileExists{microtype.sty}{% use microtype if available
  \usepackage[]{microtype}
  \UseMicrotypeSet[protrusion]{basicmath} % disable protrusion for tt fonts
}{}
\makeatletter
\@ifundefined{KOMAClassName}{% if non-KOMA class
  \IfFileExists{parskip.sty}{%
    \usepackage{parskip}
  }{% else
    \setlength{\parindent}{0pt}
    \setlength{\parskip}{6pt plus 2pt minus 1pt}}
}{% if KOMA class
  \KOMAoptions{parskip=half}}
\makeatother
\usepackage{xcolor}
\usepackage{color}
\usepackage{fancyvrb}
\newcommand{\VerbBar}{|}
\newcommand{\VERB}{\Verb[commandchars=\\\{\}]}
\DefineVerbatimEnvironment{Highlighting}{Verbatim}{commandchars=\\\{\}}
% Add ',fontsize=\small' for more characters per line
\usepackage{framed}
\definecolor{shadecolor}{RGB}{248,248,248}
\newenvironment{Shaded}{\begin{snugshade}}{\end{snugshade}}
\newcommand{\AlertTok}[1]{\textcolor[rgb]{0.94,0.16,0.16}{#1}}
\newcommand{\AnnotationTok}[1]{\textcolor[rgb]{0.56,0.35,0.01}{\textbf{\textit{#1}}}}
\newcommand{\AttributeTok}[1]{\textcolor[rgb]{0.13,0.29,0.53}{#1}}
\newcommand{\BaseNTok}[1]{\textcolor[rgb]{0.00,0.00,0.81}{#1}}
\newcommand{\BuiltInTok}[1]{#1}
\newcommand{\CharTok}[1]{\textcolor[rgb]{0.31,0.60,0.02}{#1}}
\newcommand{\CommentTok}[1]{\textcolor[rgb]{0.56,0.35,0.01}{\textit{#1}}}
\newcommand{\CommentVarTok}[1]{\textcolor[rgb]{0.56,0.35,0.01}{\textbf{\textit{#1}}}}
\newcommand{\ConstantTok}[1]{\textcolor[rgb]{0.56,0.35,0.01}{#1}}
\newcommand{\ControlFlowTok}[1]{\textcolor[rgb]{0.13,0.29,0.53}{\textbf{#1}}}
\newcommand{\DataTypeTok}[1]{\textcolor[rgb]{0.13,0.29,0.53}{#1}}
\newcommand{\DecValTok}[1]{\textcolor[rgb]{0.00,0.00,0.81}{#1}}
\newcommand{\DocumentationTok}[1]{\textcolor[rgb]{0.56,0.35,0.01}{\textbf{\textit{#1}}}}
\newcommand{\ErrorTok}[1]{\textcolor[rgb]{0.64,0.00,0.00}{\textbf{#1}}}
\newcommand{\ExtensionTok}[1]{#1}
\newcommand{\FloatTok}[1]{\textcolor[rgb]{0.00,0.00,0.81}{#1}}
\newcommand{\FunctionTok}[1]{\textcolor[rgb]{0.13,0.29,0.53}{\textbf{#1}}}
\newcommand{\ImportTok}[1]{#1}
\newcommand{\InformationTok}[1]{\textcolor[rgb]{0.56,0.35,0.01}{\textbf{\textit{#1}}}}
\newcommand{\KeywordTok}[1]{\textcolor[rgb]{0.13,0.29,0.53}{\textbf{#1}}}
\newcommand{\NormalTok}[1]{#1}
\newcommand{\OperatorTok}[1]{\textcolor[rgb]{0.81,0.36,0.00}{\textbf{#1}}}
\newcommand{\OtherTok}[1]{\textcolor[rgb]{0.56,0.35,0.01}{#1}}
\newcommand{\PreprocessorTok}[1]{\textcolor[rgb]{0.56,0.35,0.01}{\textit{#1}}}
\newcommand{\RegionMarkerTok}[1]{#1}
\newcommand{\SpecialCharTok}[1]{\textcolor[rgb]{0.81,0.36,0.00}{\textbf{#1}}}
\newcommand{\SpecialStringTok}[1]{\textcolor[rgb]{0.31,0.60,0.02}{#1}}
\newcommand{\StringTok}[1]{\textcolor[rgb]{0.31,0.60,0.02}{#1}}
\newcommand{\VariableTok}[1]{\textcolor[rgb]{0.00,0.00,0.00}{#1}}
\newcommand{\VerbatimStringTok}[1]{\textcolor[rgb]{0.31,0.60,0.02}{#1}}
\newcommand{\WarningTok}[1]{\textcolor[rgb]{0.56,0.35,0.01}{\textbf{\textit{#1}}}}
\usepackage{longtable,booktabs,array}
\usepackage{calc} % for calculating minipage widths
% Correct order of tables after \paragraph or \subparagraph
\usepackage{etoolbox}
\makeatletter
\patchcmd\longtable{\par}{\if@noskipsec\mbox{}\fi\par}{}{}
\makeatother
% Allow footnotes in longtable head/foot
\IfFileExists{footnotehyper.sty}{\usepackage{footnotehyper}}{\usepackage{footnote}}
\makesavenoteenv{longtable}
\usepackage{graphicx}
\makeatletter
\newsavebox\pandoc@box
\newcommand*\pandocbounded[1]{% scales image to fit in text height/width
  \sbox\pandoc@box{#1}%
  \Gscale@div\@tempa{\textheight}{\dimexpr\ht\pandoc@box+\dp\pandoc@box\relax}%
  \Gscale@div\@tempb{\linewidth}{\wd\pandoc@box}%
  \ifdim\@tempb\p@<\@tempa\p@\let\@tempa\@tempb\fi% select the smaller of both
  \ifdim\@tempa\p@<\p@\scalebox{\@tempa}{\usebox\pandoc@box}%
  \else\usebox{\pandoc@box}%
  \fi%
}
% Set default figure placement to htbp
\def\fps@figure{htbp}
\makeatother
\setlength{\emergencystretch}{3em} % prevent overfull lines
\providecommand{\tightlist}{%
  \setlength{\itemsep}{0pt}\setlength{\parskip}{0pt}}
\setcounter{secnumdepth}{5}
\usepackage{booktabs}
\usepackage[]{natbib}
\bibliographystyle{plainnat}
\usepackage{bookmark}
\IfFileExists{xurl.sty}{\usepackage{xurl}}{} % add URL line breaks if available
\urlstyle{same}
\hypersetup{
  pdftitle={Lie Groups and Lie Algbera},
  pdfauthor={Ashan Jayamal},
  hidelinks,
  pdfcreator={LaTeX via pandoc}}

\title{Lie Groups and Lie Algbera}
\author{Ashan Jayamal}
\date{2025-12-10}

\usepackage{amsthm}
\newtheorem{theorem}{Theorem}[chapter]
\newtheorem{lemma}{Lemma}[chapter]
\newtheorem{corollary}{Corollary}[chapter]
\newtheorem{proposition}{Proposition}[chapter]
\newtheorem{conjecture}{Conjecture}[chapter]
\theoremstyle{definition}
\newtheorem{definition}{Definition}[chapter]
\theoremstyle{definition}
\newtheorem{example}{Example}[chapter]
\theoremstyle{definition}
\newtheorem{exercise}{Exercise}[chapter]
\theoremstyle{definition}
\newtheorem{hypothesis}{Hypothesis}[chapter]
\theoremstyle{remark}
\newtheorem*{remark}{Remark}
\newtheorem*{solution}{Solution}
\begin{document}
\maketitle

{
\setcounter{tocdepth}{1}
\tableofcontents
}
\chapter{Notation and Some definitions}\label{notation-and-some-definitions}

\begin{enumerate}
\def\labelenumi{\arabic{enumi}.}
\tightlist
\item
  \(\mathrm{GL}_n(\mathbb{C})\): the group of all \(n \times n\) invertible matrices with entries in the field \(\mathbb{F}\).
\item
  \(M_n(\mathbb{C})\): the set of all \(n \times n\) matrices with entries in the field \(\mathbb{F}\).
\item
  \(U(n)\): \(n \times n\) unitary group
\item
  \(SU\): Special Unitary Group (Unitary Group + det=1)
\item
  \(O(n)\): \(n \times n\) Orthogonal group
\item
  \(SO(n)\): \(n \times n\) Special orthogonal Group (Orthogonal Group + det=1)
\end{enumerate}

\chapter{Lie Group}\label{lie-group}

\begin{definition}[Lie Group]
\protect\hypertarget{def:unnamed-chunk-1}{}\label{def:unnamed-chunk-1}A \emph{Lie group} is a smooth manifold \(G\) which is also a group, such that the group multiplication
\[\mu: G \times G \to G, \quad (g,h) \mapsto gh\]
and the inverse map
\[\iota: G \to G, \quad g \mapsto g^{-1}\]
are both smooth.
\end{definition}

\begin{example}
\protect\hypertarget{exm:unnamed-chunk-2}{}\label{exm:unnamed-chunk-2}Let \(G = \mathbb{R} \times \mathbb{R} \times S\), equipped with the group operation
\[(x_1, y_1, \lambda_1) \cdot (x_2, y_2, \lambda_2) = (x_1 + x_2,\, y_1 + y_2,\, e^{ix_1y_2}\lambda_1 \lambda_2).\]
Then \(G\) is a Lie group.
\end{example}

\begin{definition}
\protect\hypertarget{def:unnamed-chunk-3}{}\label{def:unnamed-chunk-3}Let \(A_m\) be a sequence of complex matrices in \(\mathbb{M}_n(\mathbb{C})\). We say that \(A_m \to A\) (i.e., \(A_m\) converges to a matrix \(A\)) if each entry of \(A_m\) converges to the corresponding entry of \(A\) as \(m \to \infty\); In other words,
\[(A_m)_{jk} \to A_{jk} \quad \text{for all } 1 \leq j, k \leq n.\]
\end{definition}

\begin{definition}[Matrix Lie Group]
\protect\hypertarget{def:unnamed-chunk-4}{}\label{def:unnamed-chunk-4}A \emph{matrix Lie group} is a subgroup \(G \subseteq \mathrm{GL}(n, \mathbb{C})\) with the following property: If \(\{A_m\}\) is any sequence of matrices in \(G\) such that \(A_m \to A\), then either \(A \in G\) or \(A\) is not invertible.
\end{definition}

\begin{example}[A Subgroup of $\mathrm{GL}(2, \mathbb{C})$ That Is Not a Matrix Lie Group]
\protect\hypertarget{exm:unnamed-chunk-5}{}\label{exm:unnamed-chunk-5}Let \(\alpha \in \mathbb{R} \setminus \mathbb{Q}\), and define the subgroup \(G \subseteq \mathrm{GL}(2, \mathbb{C})\) consisting of matrices of the form

\[G = \left\{
\begin{pmatrix}
e^{it} & 0 \\
0 & e^{i\alpha t}
\end{pmatrix}
: t \in \mathbb{R}
\right\}.\]

Clearly, \(G\) is a subgroup of \(\mathrm{GL}(2, \mathbb{C})\). The closure of \(G\) in the standard topology is the torus subgroup

\[\overline{G} = \left\{
\begin{pmatrix}
e^{i\theta} & 0 \\
0 & e^{i\phi}
\end{pmatrix}
: \theta, \phi \in \mathbb{R}
\right\}
= \left\{
\begin{pmatrix}
z_1 & 0 \\
0 & z_2
\end{pmatrix}
: |z_1| = |z_2| = 1
\right\},\]

which is a compact Lie group. Since \(G\) is not closed, it fails the closure condition required for matrix Lie groups. Thus, \(G\) is not a matrix Lie group.

This subgroup \(G\) is sometimes referred to as an \emph{irrational line in a torus}.
\end{example}

\begin{example}[Example of Matrix Linear Groups]
\protect\hypertarget{exm:unnamed-chunk-6}{}\label{exm:unnamed-chunk-6}\leavevmode

\begin{enumerate}
\def\labelenumi{\arabic{enumi}.}
\tightlist
\item
\end{enumerate}

\end{example}

\begin{definition}[Lie Group homomorphism and Isomorphisms]
\protect\hypertarget{def:unnamed-chunk-7}{}\label{def:unnamed-chunk-7}Let \(G\) and \(H\) be matrix Lie groups. A map \(\varphi : G \to H\) is called a \emph{Lie group homomorphism} if

\begin{enumerate}
\def\labelenumi{\roman{enumi}.}
\tightlist
\item
  \(\varphi\) is a group homomorphism, and
\item
  \(\varphi\) is continuous.
\end{enumerate}

If, in addition, \(\varphi\) is one-to-one and onto, and the inverse map \(\varphi^{-1}\) is continuous, then \(\varphi\) is called a \emph{Lie group isomorphism}.
\end{definition}

\begin{example}[Examples of Lie Group Homomorphisms]
\protect\hypertarget{exm:unnamed-chunk-8}{}\label{exm:unnamed-chunk-8}\leavevmode

\begin{enumerate}
\def\labelenumi{\arabic{enumi}.}
\item
  The determinant is a homomorphism \(\det : \mathrm{GL}_n(\mathbb{C}) \to \mathbb{C}\setminus \{0\}\)
\item
  The map \(\varphi : \mathbb{R} \to \mathrm{SO}(2)\), defined by
  \[\varphi(t) = \begin{pmatrix} \cos t & -\sin t \\ \sin t & \cos t \end{pmatrix}\]
  is also a group homomorphism.
\end{enumerate}

\end{example}

\chapter{Lie Algebra}\label{lie-algebra}

\begin{definition}[Lie Algebra]
\protect\hypertarget{def:unnamed-chunk-9}{}\label{def:unnamed-chunk-9}Let \(\mathbb{F}\) be a field. A \emph{Lie algebra} over \(\mathbb{F}\) is a vector space \(\mathfrak{g}\) over \(\mathbb{F}\) together with a map \([\cdot,\cdot] : \mathfrak{g} \times \mathfrak{g} \to \mathfrak{g}\), called the \emph{Lie bracket}, satisfying the following properties:

\begin{enumerate}
\def\labelenumi{\arabic{enumi}.}
\item
  \textbf{Bi-linearity:} The bracket \([\cdot,\cdot]\) is bilinear.
\item
  \textbf{Skew-symmetry:} For all \(X,Y \in \mathfrak{g}\),
  \[[X,Y] = -[Y,X].\]
\item
  \textbf{Jacobi identity:} For all \(X,Y,Z \in \mathfrak{g}\),
  \[[X,[Y,Z]] + [Y,[Z,X]] + [Z,[X,Y]] = 0.\]
\end{enumerate}

Two elements \(X,Y \in \mathfrak{g}\) are said to \emph{commute} if \([X,Y] = 0\). The Lie algebra \(\mathfrak{g}\) is called \emph{commutative} (or \emph{abelian}) if \([X,Y] = 0\) for all \(X,Y \in \mathfrak{g}\).
\end{definition}

\begin{example}
\protect\hypertarget{exm:unnamed-chunk-10}{}\label{exm:unnamed-chunk-10}Let \(\mathfrak{g} = \mathbb{R}^3\) and define the bracket
\[[x,y] = x \times y, \quad x,y \in \mathbb{R}^3,\]
where \(x \times y\) denotes the cross product. Then \(\mathfrak{g}\) is a Lie algebra.
\end{example}

\begin{proof}
Bilinearity and skew-symmetry are standard properties of the cross product. To verify the Jacobi identity, it suffices (by bilinearity) to check it when
\[x = e_j, \quad y = e_k, \quad z = e_\ell,\]
where \(e_1, e_2, e_3\) are the standard basis vectors of \(\mathbb{R}^3\).

\begin{itemize}
\tightlist
\item
  If \(j,k,\ell\) are all equal, each term in the Jacobi identity is zero.
\item
  If \(j,k,\ell\) are all distinct, the cross product of any two of \(e_j, e_k, e_\ell\) is a multiple of the third, so again each term in the Jacobi identity is zero.
\item
  It remains to consider the case in which two of \(j,k,\ell\) are equal and the third is different. By reordering the terms in the Jacobi identity as necessary, it suffices to verify
  \[[e_j,[e_j,e_k]] + [e_j,[e_k,e_j]] + [e_k,[e_j,e_j]] = 0. \qquad (3.1)\]
\end{itemize}

The first two terms in (3.1) are negatives of each other, and the third is zero.

Thus the Jacobi identity holds, and \(\mathfrak{g}\) is a Lie algebra.
\end{proof}

\begin{example}
\protect\hypertarget{exm:unnamed-chunk-12}{}\label{exm:unnamed-chunk-12}Let \(A\) be an associative algebra and let \(\mathfrak{g}\) be a subspace of \(A\) such that
\[XY - YX \in \mathfrak{g} \quad \text{for all } X,Y \in \mathfrak{g}.\]

Then \(\mathfrak{g}\) is a Lie algebra with bracket operation defined by
\[[X,Y] = XY - YX.\]
\end{example}

\begin{proof}
It is very easy to show the bilinear and skew symmetry properties. So, just go for Jacobi Identity.

Let \(X,Y,Z \in \mathfrak{g}\), consider

\begin{align*}
[X,[Y,Z]] &= X(YZ - ZY) - (YZ - ZY)X = XYZ - XZY - YZX + ZYX,\\
[Y,[Z,X]] &= Y(ZX - XZ) - (ZX - XZ)Y = YZX - YXZ - ZXY + XZY,\\
[Z,[X,Y]] &= Z(XY - YX) - (XY - YX)Z = ZXY - ZYX - XYZ + YXZ.
\end{align*}

Adding these three expressions together, every monomial in \(X,Y,Z\) appears exactly twice, once with a plus sign and once with a minus sign:

\begin{align*}
(XYZ - XYZ) + (XZY - XZY) + (YZX - YZX) + (YXZ - YXZ) \\+ (ZXY - ZXY) + (ZYX - ZYX) = 0.
\end{align*}
\end{proof}

\chapter{Matrix Exponential}\label{matrix-exponential}

The exponential of a matrix plays a crucial role in the theory of Lie groups. The exponential enters into the definition of the Lie algebra of a matrix Lie group and is the mechanism for passing information from the Lie algebra to the Lie group.

\begin{definition}[Matrix Exponential]
\protect\hypertarget{def:unnamed-chunk-14}{}\label{def:unnamed-chunk-14}If \(X\) is an \(n \times n\) matrix, we define the exponential of \(X\), denoted \(e^X\) or \(\exp(X)\), by the usual power series
\begin{equation}
e^X = \sum_{m=0}^{\infty} \frac{X^m}{m!}, \label{eq:matrixexp}
\end{equation}
where \(X^0=I\), and \(X^m\) denotes the repeated matrix product of \(X\) with itself.
\end{definition}

\begin{proposition}
\protect\hypertarget{prp:unnamed-chunk-15}{}\label{prp:unnamed-chunk-15}The series \eqref{eq:matrixexp} converges for all \(X \in M_n(\mathbb{C})\), and \(e^X\) is a continuous function of \(X\).
\end{proposition}

\begin{theorem}
\protect\hypertarget{thm:thm1}{}\label{thm:thm1}

Let \(X,Y \in M_n(\mathbb{C})\). Then the following properties hold:

\begin{enumerate}
\def\labelenumi{\arabic{enumi}.}
\tightlist
\item
  \(e^0 = I\).
\item
  \(e^{X^T}=\left( e^X\right)^T\), \((e^X)^* = e^{X^*}\).
\item
  \(e^X\) is invertible and \((e^X)^{-1} = e^{-X}\).
\item
  If \(A\) is invertible \(n \times n\) matrix, then \(e^{AXA^{-1}}=Ae^XA^{-1}\)
\item
  \(e^{(\alpha+\beta)X} = e^{\alpha X} e^{\beta X}\) for all \(\alpha,\beta \in \mathbb{C}\).
\item
  If \(XY = YX\), then \(e^{X+Y} = e^X e^Y = e^Y e^X\).
\item
  If \(C \in GL_n(\mathbb{C})\), then \(e^{CXC^{-1}} = C e^X C^{-1}\).
\end{enumerate}

\end{theorem}

\begin{proof}
Point (1) is immediate from the definition of the exponential series. Point (2) follows by taking term-by-term adjoints in the series for \(e^X\). Points (3) and (4) are special cases of Point (5).

To verify Point (5), note that both series \(e^X\) and \(e^Y\) converge absolutely, so we may multiply them term by term:

\begin{equation}
e^X e^Y = \sum_{m=0}^{\infty} \sum_{k=0}^{m} \frac{X^k}{k!} \frac{Y^{\,m-k}}{(m-k)!}
= \sum_{m=0}^{\infty} \frac{1}{m!} \sum_{k=0}^{m} \binom{m}{k} X^k Y^{\,m-k}. \label{eq:exp25}
\end{equation}

If \(X\) and \(Y\) commute, then
\[(X+Y)^m = \sum_{k=0}^{m} \binom{m}{k} X^k Y^{\,m-k},\]
so \eqref{eq:exp25} becomes
\[e^X e^Y = \sum_{m=0}^{\infty} \frac{(X+Y)^m}{m!} = e^{X+Y}.\]

For Point (6), observe that
\[(CXC^{-1})^m = C X^m C^{-1},\]
so term by term we have
\[e^{CXC^{-1}} = \sum_{m=0}^{\infty} \frac{(CXC^{-1})^m}{m!}
= \sum_{m=0}^{\infty} \frac{C X^m C^{-1}}{m!}
= C \left( \sum_{m=0}^{\infty} \frac{X^m}{m!} \right) C^{-1}
= C e^X C^{-1}.\]
\end{proof}

\section{Computing Exponential}\label{computing-exponential}

If \(X\) is diagonalizable, then there exist invertible \(A\) and \(\lambda_1,...,\lambda_n\) such that

\[X = A
\begin{bmatrix}
{\lambda_1} & 0 & \cdots & 0 \\
0 & {\lambda_2} & \cdots & 0 \\
\vdots & \vdots & \ddots & \vdots \\
0 & 0 & \cdots & {\lambda_n}
\end{bmatrix}
A^{-1}\]

Then by Theorem \ref{thm:thm1} fourth property, we get the following:

\[e^X = A
\begin{bmatrix}
e^{\lambda_1} & 0 & \cdots & 0 \\
0 & e^{\lambda_2} & \cdots & 0 \\
\vdots & \vdots & \ddots & \vdots \\
0 & 0 & \cdots & e^{\lambda_n}
\end{bmatrix}A^{-1}\]

If \(X\) is nilpotent then the series that defines \(e^X\) terminates.

\begin{example}
\protect\hypertarget{exm:unnamed-chunk-17}{}\label{exm:unnamed-chunk-17}Consider the matrices

\[X_1 =
\begin{bmatrix}
0 & -a \\
a & 0
\end{bmatrix}, \quad
X_2 =
\begin{bmatrix}
0 & a & b \\
0 & 0 & c \\
0 & 0 & 0
\end{bmatrix}, \quad
X_3 =
\begin{bmatrix}
a & b \\
0 & a
\end{bmatrix}.\]

Then

\[e^{X_1} =
\begin{bmatrix}
\cos a & -\sin a \\
\sin a & \cos a
\end{bmatrix}, \text{ and }
e^{X_2} =
\begin{bmatrix}
1 & a & b + \tfrac{ac}{2} \\
0 & 1 & c \\
0 & 0 & 1
\end{bmatrix}, \text{ and } 
e^{X_3} =
\begin{bmatrix}
e^a & e^a b \\
0 & e^a
\end{bmatrix}.\]
\end{example}

\begin{proof}
\[e^{X_1} = 
\begin{bmatrix} 1 & i \\ i & 1 \end{bmatrix}
\begin{bmatrix} e^{-ia} & 0 \\ 0 & e^{ia} \end{bmatrix}
\frac{1}{2}
\begin{bmatrix} 1 & -i \\ -i & 1 \end{bmatrix}.\]
\end{proof}

\chapter{The Lie Algebra of a Matrix Lie Group}\label{the-lie-algebra-of-a-matrix-lie-group}

Now we going to discuss the Lie Algebra \(\mathfrak{g}\) that associated with Matrix Lie group \(G\).

\begin{definition}
\protect\hypertarget{def:unnamed-chunk-19}{}\label{def:unnamed-chunk-19}If \(G \subset GL(n, \mathbb{C})\) is a matrix Lie group, then the Lie algebra \(\mathfrak{g}\) of \(G\) is defined as follows:
\[
\mathfrak{g} = \left\{ X \in M_n(\mathbb{C}) \;\middle|\; e^{tX} \in G \;\; \text{for all } t \in \mathbb{R} \right\}.
\]
\end{definition}

\begin{proposition}
\protect\hypertarget{prp:unnamed-chunk-20}{}\label{prp:unnamed-chunk-20}

For any matrix Lie group \(G\), the Lie algebra \(\mathfrak{g}\) of \(G\) has the following properties:

\begin{itemize}
\tightlist
\item
  The zero matrix \(0\) belongs to \(\mathfrak{g}\).
\item
  For all \(X \in \mathfrak{g}\), \(tX \in \mathfrak{g}\) for all real numbers \(t\).
\item
  For all \(X, Y \in \mathfrak{g}\), \(X + Y \in \mathfrak{g}\).
\item
  For all \(A \in G\) and \(X \in \mathfrak{g}\) we have \(AXA^{-1} \in \mathfrak{g}\).
\item
  For all \(X, Y \in \mathfrak{g}\), the commutator \([X, Y] := XY - YX\) belongs to \(\mathfrak{g}\).
\end{itemize}

\end{proposition}

\begin{proof}
Write the proof here
\end{proof}

\begin{remark}
Write about real vector spaces.
\end{remark}

\section{Later updared with the correct picture.}\label{later-updared-with-the-correct-picture.}

\chapter{Blocks}\label{blocks}

\section{Equations}\label{equations}

Here is an equation.

\begin{equation} 
  f\left(k\right) = \binom{n}{k} p^k\left(1-p\right)^{n-k}
  \label{eq:binom}
\end{equation}

You may refer to using \texttt{\textbackslash{}@ref(eq:binom)}, like see Equation \eqref{eq:binom}.

\section{Theorems and proofs}\label{theorems-and-proofs}

Labeled theorems can be referenced in text using \texttt{\textbackslash{}@ref(thm:tri)}, for example, check out this smart theorem \ref{thm:tri}.

\begin{theorem}
\protect\hypertarget{thm:tri}{}\label{thm:tri}For a right triangle, if \(c\) denotes the \emph{length} of the hypotenuse
and \(a\) and \(b\) denote the lengths of the \textbf{other} two sides, we have
\[a^2 + b^2 = c^2\]
\end{theorem}

Read more here \url{https://bookdown.org/yihui/bookdown/markdown-extensions-by-bookdown.html}.

\section{Callout blocks}\label{callout-blocks}

The R Markdown Cookbook provides more help on how to use custom blocks to design your own callouts: \url{https://bookdown.org/yihui/rmarkdown-cookbook/custom-blocks.html}

\chapter{Sharing your book}\label{sharing-your-book}

\section{Publishing}\label{publishing}

HTML books can be published online, see: \url{https://bookdown.org/yihui/bookdown/publishing.html}

\section{404 pages}\label{pages}

By default, users will be directed to a 404 page if they try to access a webpage that cannot be found. If you'd like to customize your 404 page instead of using the default, you may add either a \texttt{\_404.Rmd} or \texttt{\_404.md} file to your project root and use code and/or Markdown syntax.

\section{Metadata for sharing}\label{metadata-for-sharing}

Bookdown HTML books will provide HTML metadata for social sharing on platforms like Twitter, Facebook, and LinkedIn, using information you provide in the \texttt{index.Rmd} YAML. To setup, set the \texttt{url} for your book and the path to your \texttt{cover-image} file. Your book's \texttt{title} and \texttt{description} are also used.

This \texttt{gitbook} uses the same social sharing data across all chapters in your book- all links shared will look the same.

Specify your book's source repository on GitHub using the \texttt{edit} key under the configuration options in the \texttt{\_output.yml} file, which allows users to suggest an edit by linking to a chapter's source file.

Read more about the features of this output format here:

\url{https://pkgs.rstudio.com/bookdown/reference/gitbook.html}

Or use:

\begin{Shaded}
\begin{Highlighting}[]
\NormalTok{?bookdown}\SpecialCharTok{::}\NormalTok{gitbook}
\end{Highlighting}
\end{Shaded}


  \bibliography{book.bib,packages.bib}

\end{document}
